\origpage{113}{GRAPH MINORS. XV. GIANT STEPS}
\noindent closed curve in the surface that does not bound a disc meets the drawing many times). Our initial pair $H_0, \Sigma_0$ is a degenerate case, because for the sphere $\Sigma_0$ representativity does not make sense, but we showed in earlier papers that in the sphere, high tree-width is the appropriate analogue for high representativity. Let us make a sequence $H_i, \Sigma_i$ ($i = 0, 1, 2, ...$) of such pairs, at each stage sacrificing a large amount of the remaining representativity if necessary for an increase in genus (we permit $\Sigma_i$ to be non-orientable, in which case its ``genus'' is half the number of crosscaps). Since $H_0, \Sigma_0$ has large representativity, the initial few terms in our sequence still have high representativity. Now every graph with high representativity that can be drawn in a surface in which $K_n$ can be drawn has a $K_n$ minor. Hence the process stops in the initial few terms, before $\Sigma_i$ has grown complex enough that $K_n$ can be drawn in it. We therefore have a pair $H, \Sigma$ (say), where $\Sigma$ has bounded genus and $H$ is drawn in it with very large representativity, so that no improvement in $\Sigma$ can be made.

How does the remainder of $G$ attach to the subgraph $H$? The theorem that we are going to prove here says roughly that either

\hspace*{12pt}(i) we can delete a small area of $H$, and replace it by some other subgraph of $G$, to obtain a pair $H', \Sigma'$ where $\Sigma'$ is obtained by adding a crosscap to $\Sigma$ and $H'$ still has large representativity--but this contradicts the choice of $H, \Sigma$;

\hspace*{12pt}(ii) we can choose another $H', \Sigma$ so that $H'$ can be drawn in $\Sigma$ with crossings, but with a large constant number of crossings, all pairwise ``far apart'' and still with high representativity in the appropriate sense--but this implies that $G$ has a $K_n$ minor, a contradiction;

\hspace*{12pt}(iii) we can rearrange a small area of $H$ so that some path of $G$ has both ends far apart in $H$ and is otherwise disjoint from $H$; or,

\hspace*{12pt}(iv) $G$ has the desired structure (roughly, all of $G$ can be drawn in $\Sigma$ ``up to 3-separations,'' except for a bounded number of local areas of non-planarity, called ``vortices'').

Thus, we are done in every case except (iii). That case needs work, and to handle it we actually need a more complicated optimization of $H, \Sigma$ than is described here; but we hope this provides some motivation for the present result.

This paper is closely related to [5], and uses most of the same terminology. In particular, the reader should see [5, Sections 2--4] for the meaning of the following terms and notation: $G_1 \cup G_2$, $G_1 \cap G_2$, separation, order of a separation, tangle, order of a tangle, free, open and closed disc, surface, $bd(\Sigma)$, $O$-arc, line, ends of a line, drawing, $U(H)$, region, $R(H)$, atom, $A(H)$, radical drawing, respectful tangle, metric of a tangle, $\lambda$-zone, clearing, 2-cell, rigid, dial, natural order, society, $\bar{\Omega}$, segregation, 3-segregation,
